\subsection{CONTINUOUS-POLYNOMIALS WITH VALUES IN g}

Let I be a finite set and let \(\mathbf{P}(\mathfrak{g}^{\mathrm{I}};\mathfrak{g})\) (resp. \(\hat{\mathbf{P}} (\mathfrak{g}^{\mathrm{I}};\mathfrak{g}))\) be the vector space of continuous-polynomials (resp. formal power series with continuous components) on \(\mathbf{g}^{\mathrm{I}}\) with values in \(\mathfrak{g}\) . Recall (Differentiable and Analytic Manifolds, R, Appendix) that \(\mathbf{P}(\mathfrak{g}^{\mathrm{I}};\mathfrak{g})\) has a graduation of type \(\mathbf{N}^{\mathrm{I}}\) and that \(\hat{\mathbf{P}} (\mathfrak{g}^{\mathrm{I}};\mathfrak{g})\) is identified with the completion of the vector space \(\mathbf{P} (\mathfrak{g}^{\mathrm{I}};\mathfrak{g})\) with the topology defined by the filtration associated with the graduation of \(\mathbf{P}(\mathfrak{g}^{\mathrm{I}};\mathfrak{g})\) . Moreover, \(\mathbf{P}(\mathfrak{g}^{\mathrm{I}};\mathfrak{g})\) is a graded Lie algebra with the bracket defined by \([f,g](\pmb {x}) = [f(\pmb {x}),g(\pmb {x})]\) for \(f, g\) in \(\mathrm{P}(\mathfrak{g}^{\mathrm{I}}; \mathfrak{g}), x \in \mathfrak{g}^{\mathrm{I}}\) ; this Lie algebra structure can be extended by continuity to \(\hat{\mathrm{P}}(\mathfrak{g}^{\mathrm{I}}; \mathfrak{g})\) and makes it into a complete Hausdorff filtered Lie algebra.

By proposition 2 of § 6, no. 3, there exists one and only one continuous Lie algebra homomorphism \(\phi_{\mathrm{I}}:u\mapsto \tilde{u}\) of \(\hat{\mathbf{L}} (\mathbf{I})\) into \(\hat{\mathbf{P}} (\mathbf{g}^{\mathrm{I}};\mathfrak{g})\) mapping the indeterminate of index \(i\) to \(\mathsf{pr}_{\mathrm{i}}\) for all \(i\in \mathbb{I}\) , since \(\mathsf{pr}_{\mathrm{i}}\in \mathsf{P}(\mathfrak{g}^{\mathrm{I}};\mathfrak{g})\) . It follows that \(\tilde{u}\in \mathsf{P}(\mathfrak{g}^{\mathrm{I}};\mathfrak{g})\) for \(u\in \mathbf{L}(\mathbf{I})\) ; more precisely, when \(u\in \mathbf{L}(\mathbf{I})\) , \(\tilde{u}\) is just the polynomial mapping \((t_i)\mapsto u((t_i))\) of § 2, no. 4. On the other hand, clearly \(\phi_{\mathrm{I}}\) is compatible with the multigraduations of \(\mathbf{L}(\mathbf{I})\) and \(\mathbf{P}(\mathfrak{g}^{\mathrm{I}};\mathfrak{g})\) . If \(u = \sum_{\nu \in \mathbf{N}^{\mathrm{I}}}u_{\nu}\) , where \(u_{\nu}\in \mathbf{L}^{\nu}(\mathbf{I})\) for \(\nu \in \mathbf{N}^{\mathrm{I}}\) , then

\[
\tilde {u} = \sum_ {\nu \in \mathbf {N} ^ {\mathrm{I}}}  \tilde {u} _ {\nu}, \quad \text { where } \tilde {u} _ {\nu} \in \mathrm{P} _ {\nu} (\mathfrak {g} ^ {\mathrm{I}}; \mathfrak {g}).
\]

Let \(\pmb{u} = (u_j)_{j\in J}\) be a finite family of elements of \(\hat{\mathbf{L}} (\mathbf{I})\) , let \(v\in \hat{\mathbf{L}} (\mathbf{J})\) and let \(w = v\circ \pmb {u}\) (§ 6, no. 3). We write \(\tilde{\pmb{u}} = (\tilde{u}_j)_{j\in \mathbb{J}}\in \mathfrak{J}\) . Then

\[
\tilde {v} \circ \tilde {\boldsymbol {u}} = (v \circ \boldsymbol {u}) ^ {\sim}.\tag{2}
\]

This follows by extending by continuity formula (7) of §6, no. 3 and from Differentiable and Analytic Manifolds, R, Appendix, no. 6.

